\documentclass[11pt,a4paper]{article}
\usepackage[T1]{fontenc}
\usepackage[utf8]{inputenc}
\usepackage{charter}
\usepackage[scaled=0.86]{helvet}
\usepackage{amsmath,amssymb}
\usepackage{booktabs}
\usepackage{array}
\usepackage{xcolor}
\usepackage{microtype}
\usepackage[margin=2.3cm]{geometry}
\usepackage[colorlinks=true,linkcolor=blue!50!black,urlcolor=blue!55!black]{hyperref}
\usepackage{titlesec}
\usepackage{tcolorbox}

\definecolor{bleufonce}{RGB}{31,78,156}
\definecolor{rougefonce}{RGB}{155,34,38}
\definecolor{grisclair}{RGB}{245,246,248}

\titleformat{\section}{\large\bfseries\color{bleufonce}}{\thesection}{0.7em}{}
\titleformat{\subsection}{\normalsize\bfseries}{\thesubsection}{0.6em}{}
\titlespacing*{\section}{0pt}{13pt}{5pt}
\titlespacing*{\subsection}{0pt}{9pt}{3pt}

\newtcolorbox{encadre}{colback=grisclair,colframe=bleufonce,boxrule=0.5pt,
  left=6pt,right=6pt,top=5pt,bottom=5pt}
\newtcolorbox{choix}{colback=white,colframe=rougefonce,boxrule=0.5pt,
  left=6pt,right=6pt,top=5pt,bottom=5pt}

\setlength{\parskip}{4pt}
\setlength{\parindent}{0pt}
\frenchspacing

\newcommand{\ou}[1]{\textit{\small Où cela figure dans le rapport : #1}}

\title{\vspace{-1.6cm}\bfseries Note méthodologique\\
\large Valeurs de référence, erreur type, pente de convergence\\
et nombre de trajectoires}
\author{\normalsize Moteur de pricing d'options --- annexe technique}
\date{\normalsize\today}

\begin{document}
\maketitle
\vspace{-0.7cm}

\begin{encadre}
\textbf{Objet.} Le rapport annonce des prix d'exotiques, des erreurs types, des
pentes de convergence et des nombres de trajectoires. Cette note dit
\emph{comment chacun de ces nombres est obtenu} : avec quels paramètres, sur
combien de trajectoires, avec quelle graine, et par quel calcul exact. Chaque
section indique le tableau ou la section du rapport concerné.

Elle signale aussi, en encadré rouge, chaque \textbf{choix arbitraire} :
là où plusieurs options étaient défendables, laquelle a été retenue et pourquoi.
\end{encadre}

\section{Le point de référence, commun à tout le rapport}

Sauf mention contraire, tous les chiffres du rapport et de cette note portent
sur le même contrat :

\begin{center}
\small
\begin{tabular}{lll}
\toprule
Spot $S_0$ & 100 & prix courant du sous-jacent \\
Strike $K$ & 100 & à la monnaie \\
Taux sans risque $r$ & 5 \% & composition continue \\
Volatilité $\sigma$ & 20 \% & annualisée \\
Maturité $T$ & 1 an & \\
Rendement $q$ & 0 & action sans dividende \\
Fixings $n$ & \textbf{252} & pour les produits path-dependent uniquement \\
Barrière $B$ & 120 & pour les options à barrière \\
\bottomrule
\end{tabular}
\end{center}

\begin{choix}
\textbf{Choix : $n = 252$.} C'est le nombre de jours de bourse dans une année.
Ce choix n'est pas neutre et doit être cité à côté de chaque prix
d'exotique : une barrière \emph{out} observée 100 fois par an vaut 1,42, la
même observée 252 fois vaut 1,33. Moins on observe, moins l'option a
d'occasions de mourir, donc plus elle est chère. Même mécanisme sur
l'asiatique : moins de fixings, moins de lissage, prix plus élevé.
\end{choix}

\section{L'erreur type}

\ou{tableaux 6, 7, 8 et 10 --- toutes les colonnes « erreur type » et « IC95 ».}

\subsection{Ce que c'est}

Un prix Monte-Carlo est une \textbf{moyenne}. On tire $M$ trajectoires, on
calcule le gain de l'option sur chacune, on actualise, on fait la moyenne.

Relancé avec d'autres tirages, le calcul ne rend pas exactement le même prix.
\textbf{L'erreur type mesure de combien ce prix bougerait d'une exécution à
l'autre.} C'est l'incertitude propre à la méthode numérique : même un moteur
parfaitement juste la subit.

\subsection{Deux quantités qu'il ne faut pas confondre}

\begin{itemize}\setlength\itemsep{1pt}
\item \textbf{$\omega$, l'écart-type des gains.} Dispersion des trajectoires
      \emph{individuelles} : certaines rapportent zéro, d'autres beaucoup.
      C'est une propriété du produit. \textbf{Elle ne diminue pas quand on
      simule davantage.}
\item \textbf{SE, l'erreur type.} Dispersion de la \emph{moyenne}. Elle vaut
      $\omega/\sqrt{M}$, et elle diminue.
\end{itemize}

\subsection{Le calcul, vérifié à la main}

Asiatique arithmétique, $M = 50\,000$, $n = 252$, graine 12345 :

\begin{center}
\small
\begin{tabular}{lr}
\toprule
Moyenne des gains actualisés (le prix) & 5,786023 \\
Écart-type des gains $\omega$ & \textbf{8,046696} \\
Erreur type $=\ \omega/\sqrt{50\,000}$ & \textbf{0,035986} \\
Valeur rendue par le moteur & 0,035986 \\
\midrule
Part des trajectoires qui rapportent zéro & 44,4 \% \\
Demi-largeur IC 95 \% $=\ 1{,}96 \times 0{,}035986$ & 0,0705 \\
\bottomrule
\end{tabular}
\end{center}

La dispersion des gains ($8{,}05$) vaut 1,4 fois le prix lui-même. C'est
normal : 44 \% des trajectoires ne rapportent rien, une minorité rapporte
beaucoup. Diviser par $\sqrt{M}$ ramène cette dispersion à trois centimes et
demi.

\begin{choix}
\textbf{Choix : $\pm 1{,}96$ et non $\pm 2$ ou $\pm 3$.} Le facteur 1,96 donne
un intervalle à 95 \%, la convention usuelle. Il repose sur le théorème central
limite, légitime ici car $M \geq 10^4$, malgré l'asymétrie extrême de la
distribution des gains. À $M$ très petit (quelques centaines), l'intervalle
serait optimiste.
\end{choix}

\begin{encadre}
\textbf{Le piège du mode antithétique.} Avec des variables antithétiques, les
tirages vont par paires corrélées : ils ne sont plus indépendants. L'erreur
type doit alors être calculée sur les \emph{moyennes de paires}, qui sont
indépendantes entre elles, et non sur l'échantillon aplati. Appliquer
$\omega/\sqrt{M}$ à l'ensemble des tirages conduit à ne mesurer aucun gain et à
conclure à tort que la technique est inefficace. Le prix reste correct ; c'est
l'instrument de mesure de sa précision qui est faussé. Le moteur gère cela dans
\texttt{\_summarize}.
\end{encadre}

\section{Les valeurs de référence}

\subsection{Le principe : l'ancrage}

Une option européenne a une valeur exacte, donnée par Black--Scholes. Une
asiatique arithmétique ou une barrière n'en ont aucune : ni formule, ni prix
coté sur le marché. Il faut donc \emph{fabriquer} une valeur de référence, et
savoir dire à quel point on peut lui faire confiance.

Le rapport utilise quatre références. Deux sont exactes, deux sont simulées.

\begin{center}
\small
\begin{tabular}{llll}
\toprule
Référence & Valeur & Nature & Incertitude \\
\midrule
Call européen & 10,4506 & formule Black--Scholes & \textbf{aucune} \\
Asiatique géométrique & 5,5655 & formule fermée & \textbf{aucune} \\
Asiatique arithmétique & 5,7822 & simulation + contrôle & $\pm$ 0,0004 \\
Barrière up-and-in & 9,1235 & simulation + contrôle & $\pm$ 0,0077 \\
Barrière up-and-out & 1,3271 & simulation, sans contrôle & $\pm$ 0,0077 \\
\bottomrule
\end{tabular}
\end{center}

\subsection{Référence 1 --- le call européen : 10,4506}

\ou{tableaux 6, 7, 8 et 14, sections 5.4 et 11.4.}

Aucune simulation. C'est la formule de Black--Scholes appliquée au point de
référence. Cette valeur ne porte aucune incertitude : c'est un nombre calculé,
pas estimé. Elle sert d'étalon aux trois méthodes du moteur.

\subsection{Référence 2 --- l'asiatique géométrique : 5,5655}

\ou{section 8.1, tableau 14, section 11.4 niveau 3.}

Aucune simulation non plus. Un produit de variables log-normales reste
log-normal : la moyenne géométrique des 252 prix a donc une loi connue, et le
prix s'écrit sous une forme de type Black--Scholes (équation 30 du rapport).

Cette option n'a pratiquement aucun usage commercial. Elle est précieuse pour
deux raisons seulement : elle valide le pricer exotique sur un cas où la vérité
est connue, et elle sert de variable de contrôle à l'asiatique arithmétique.

\subsection{Référence 3 --- l'asiatique arithmétique : 5,7822}

\ou{tableaux 6, 8 et 14, section 4.7.}

C'est la référence la plus construite du projet. Aucune formule n'existe :
une somme de variables log-normales n'est pas log-normale.

\subsubsection*{Tous les paramètres}

\begin{center}
\small
\begin{tabular}{ll}
\toprule
Payoff & $\max(\text{moyenne arithmétique des 252 prix} - K,\ 0)$ \\
Trajectoires par lot & \textbf{100\,000} \\
Nombre de lots & \textbf{10} \\
Trajectoires au total & \textbf{1\,000\,000} \\
Tirages gaussiens au total & 252 millions \\
Graines & 999, 1000, \dots, 1008 (une par lot) \\
Variable de contrôle & asiatique géométrique, prix exact 5,565509 \\
Coefficient $\beta$ & estimé par régression sur chaque lot \\
Antithétiques & non \\
\midrule
Résultat & \textbf{5,78223} \\
Incertitude (IC 95 \%) & \textbf{$\pm$ 0,00044} \\
Précision visée par l'étude & $\pm$ 0,01 \\
Marge & \textbf{23 fois plus fine que la cible} \\
\bottomrule
\end{tabular}
\end{center}

\subsubsection*{Le calcul déroulé, sur le lot n\textsuperscript{o} 0}

\textbf{Étape 1 --- simuler.} 100\,000 trajectoires de 252 pas sous la
probabilité risque-neutre (dérive $r-q$), graine 999.

\textbf{Étape 2 --- deux gains par trajectoire.} Sur la \emph{même} trajectoire
on calcule la moyenne arithmétique des 252 prix, qui donne le gain $A$ du
produit qu'on veut valoriser, et la moyenne géométrique, qui donne le gain $G$
du contrôle. Les deux sont actualisés.

\textbf{Étape 3 --- mesurer l'erreur de la simulation sur le contrôle.} On
connaît la vraie valeur de la géométrique, donc on peut voir de combien la
simulation se trompe dessus :

\begin{center}
\small
\begin{tabular}{lr}
\toprule
Moyenne des gains géométriques simulés & 5,574205 \\
Prix exact de la géométrique (formule fermée) & 5,565509 \\
\midrule
\textbf{Erreur observée de la simulation} & \textbf{$+$0,008696} \\
\bottomrule
\end{tabular}
\end{center}

Ce lot a surestimé la géométrique de 0,0087. Comme les deux options portent sur
les mêmes trajectoires, il a très probablement surestimé l'arithmétique d'un
montant comparable.

\textbf{Étape 4 --- corriger.} On estime par régression de quel facteur $\beta$
l'erreur se transmet de $G$ vers $A$, puis on retranche :

\begin{center}
\small
\begin{tabular}{lr}
\toprule
Moyenne brute des gains arithmétiques & 5,790842 \\
$\beta$ estimé $=\ \mathrm{Cov}(A,G)/\mathrm{Var}(G)$ & 1,0350 \\
Correction $=\ \beta \times 0{,}008696$ & $-$0,009002 \\
\midrule
\textbf{Prix du lot 0} & \textbf{5,781840} \\
\bottomrule
\end{tabular}
\end{center}

\begin{encadre}
\textbf{Ce que la correction fait à la dispersion.} Avant correction,
l'écart-type des gains vaut \textbf{7,99}. Après, celui des gains corrigés
tombe à \textbf{0,22}. L'erreur type du lot passe de 0,0253 à \textbf{0,0007} :
un facteur \textbf{36}, pour le même nombre de trajectoires et un temps de
calcul quasi identique.
\end{encadre}

\textbf{Étape 5 --- répéter et moyenner.} Neuf autres lots, graines 1000 à
1008 :

\begin{center}\small
5,78184 \quad 5,78271 \quad 5,78230 \quad 5,78300 \quad 5,78153 \\
5,78309 \quad 5,78265 \quad 5,78215 \quad 5,78078 \quad 5,78223
\end{center}

\begin{center}
\small
\begin{tabular}{lr}
\toprule
\textbf{Référence} $=$ moyenne des dix & \textbf{5,78223} \\
Écart-type des dix & 0,00071 \\
Erreur type $=\ 0{,}00071/\sqrt{10}$ & 0,00022 \\
\textbf{Incertitude} $=\ 1{,}96 \times 0{,}00022$ & \textbf{$\pm$ 0,00044} \\
\bottomrule
\end{tabular}
\end{center}

\subsubsection*{L'appel de code}

\begin{center}\small
\begin{minipage}{0.93\linewidth}
\ttfamily\footnotesize
prix = [mc\_asian(S=100, K=100, r=0.05, sigma=0.20, T=1.0,\\
\phantom{prix = [mc\_asian(}n\_steps=252, n\_paths=100\_000,\\
\phantom{prix = [mc\_asian(}control\_variate=True,\\
\phantom{prix = [mc\_asian(}rng=np.random.default\_rng(999 + k)).price\\
\phantom{prix = [}for k in range(10)]\\[3pt]
reference = mean(prix)\\
incertitude = 1.96 * std(prix, ddof=1) / sqrt(10)
\end{minipage}
\end{center}

C'est exactement ce que fait \texttt{reference\_par\_lots} dans
\texttt{tools/etude\_convergence.py}.

\subsection{Référence 4 --- les barrières : 1,3271 et 9,1235}

\ou{tableaux 10 et 14, section 8.2.1.}

Recette identique : mêmes 10 lots de 100\,000, mêmes graines 999 à 1008,
$n = 252$. Seul le contrôle change.

\begin{center}
\small
\begin{tabular}{lccc}
\toprule
Produit & Contrôle & Référence & Incertitude \\
\midrule
up-and-\textbf{in}  & vanille (Black--Scholes, prix exact) & 9,1235 & $\pm$ 0,0077 \\
up-and-\textbf{out} & \emph{aucun ne fonctionne} & 1,3271 & $\pm$ 0,0077 \\
\bottomrule
\end{tabular}
\end{center}

\begin{choix}
\textbf{Limite assumée : la référence de la up-and-out.} Elle est 17 fois moins
précise que celle de l'asiatique, et seulement 1,3 fois plus fine que la cible
du centime. Ce n'est pas un défaut d'implémentation : la corrélation entre le
gain de la up-and-out et celui de la vanille vaut $-0{,}01$, aucun contrôle
n'est possible (section 4.4 de cette note). La seule façon d'améliorer cette
référence serait d'augmenter le nombre de lots. C'est le point faible du
projet, et il vaut mieux l'énoncer que le masquer.
\end{choix}

\begin{encadre}
\textbf{Attention au test de parité in $+$ out $=$ vanille.} Il donne un écart
de $-0{,}0000$, ce qui semble spectaculaire. Ce n'est pas une vérification
indépendante : sur chaque trajectoire, in $+$ out $=$ vanille par
construction, donc contrôler la « in » par la vanille avec $\beta \simeq 1$
revient algébriquement à calculer « vanille moins out ». La somme redonne donc
la vanille quoi qu'il arrive. Ce test contrôle l'implémentation, pas le
résultat.
\end{encadre}

\subsection{Pourquoi la graine des références diffère de celle des mesures}

Les références utilisent les graines 999 et suivantes, les mesures la graine
12345. Ce n'est pas cosmétique.

NumPy remplit la matrice des trajectoires ligne par ligne. Les 200\,000
premières lignes d'un tirage de 400\,000 sont donc \emph{exactement} le tirage
de taille 200\,000. Si la référence et la mesure partagent la graine, vérifier
que « la référence tombe dans l'intervalle de confiance » revient à comparer un
échantillon à un sur-ensemble de lui-même : le contrôle devient partiellement
auto-réalisateur. Changer la graine coûte un caractère et rend le test honnête.

\section{Les lots : pourquoi dix paquets plutôt qu'un seul gros calcul}

\subsection{Ce que fait la méthode}

Au lieu d'une simulation unique d'un million de trajectoires, on en lance dix
indépendantes de cent mille, avec dix graines différentes. La référence est la
moyenne des dix prix ; son incertitude est l'écart-type de ces dix prix divisé
par $\sqrt{10}$. C'est la méthode des \emph{batch means}.

\subsection{Est-ce équivalent à un million de trajectoires ? Oui}

Dix lots de cent mille, c'est un million de trajectoires, ni plus ni moins.
Aucune information n'est perdue, et la précision est la même :

\begin{center}
\small
\begin{tabular}{lc}
\toprule
Erreur type mesurée par la dispersion des dix lots & 0,000224 \\
Erreur type prédite par $\text{SE}(100\,000)/\sqrt{10}$ & 0,000222 \\
\bottomrule
\end{tabular}
\end{center}

Les deux coïncident à 1 \% près. Un unique appel à $M = 1\,000\,000$ aurait
donné le même résultat --- s'il avait pu tourner.

\subsection{Alors pourquoi le faire ? Trois raisons}

\begin{enumerate}\setlength\itemsep{2pt}
\item \textbf{La mémoire.} C'est la raison décisive. Une simulation d'un
      million de trajectoires sur 252 pas demande $10^6 \times 252 \times 8$
      octets, soit \textbf{2 Go d'un seul tenant}, et le calcul échoue à
      l'allocation. Dix lots de cent mille n'occupent jamais plus de 200 Mo à
      la fois. C'est exactement la contrainte identifiée en section 11.5 du
      rapport.
\item \textbf{Une incertitude sans hypothèse.} On \emph{observe} directement de
      combien le prix bouge d'un lot à l'autre. Les dix nombres sont sous les
      yeux.
\item \textbf{Des lots vraiment indépendants.} Chaque lot a sa propre graine :
      aucun ne partage de trajectoire avec un autre, ni avec les mesures.
\end{enumerate}

\begin{choix}
\textbf{Choix : dix lots de cent mille.} Ce n'est pas obligatoire. Si la
mémoire le permettait, un unique appel donnerait le même résultat. Le nombre 10
est un compromis : en dessous de 5 lots, l'écart-type calculé sur si peu de
points est lui-même trop bruité pour donner une incertitude fiable ; au-delà
de 20, on paie du temps de calcul pour un gain marginal. Le nombre 100\,000
est la plus grosse taille qui tient confortablement en mémoire sur 252 pas.
\end{choix}

\section{La pente de convergence}

\ou{figure 10 et sa légende, section 7.3.}

\subsection{Le calcul, sur des chiffres réels}

L'erreur type est mesurée à quatre tailles d'échantillon (asiatique
arithmétique, Monte-Carlo seul, graine 12345), puis on ajuste une droite sur
les logarithmes :

\begin{center}
\small
\begin{tabular}{rrrr}
\toprule
$M$ & erreur type & $\ln M$ & $\ln(\text{SE})$ \\
\midrule
25\,000  & 0,05070 & 10,1266 & $-$2,9819 \\
50\,000  & 0,03599 & 10,8198 & $-$3,3246 \\
100\,000 & 0,02535 & 11,5129 & $-$3,6750 \\
200\,000 & 0,01789 & 12,2061 & $-$4,0238 \\
\bottomrule
\end{tabular}
\end{center}

La régression linéaire de $\ln(\text{SE})$ sur $\ln M$ donne une droite de
coefficient directeur $-0{,}5015$. On affiche l'opposé, soit une
\textbf{pente de 0,501}.

Vérification à la main sur les deux points extrêmes seulement :
$$
\frac{-2{,}9819 - (-4{,}0238)}{12{,}2061 - 10{,}1266} = \frac{1{,}0419}{2{,}0795}
= 0{,}501 .
$$
Le même nombre. La régression sur quatre points ne fait que lisser.

\subsection{Ce que la pente prouve, et ce qu'elle ne prouve pas}

\begin{encadre}
La pente \textbf{n'est pas une performance}. Elle vaut $1/2$ pour tout
estimateur Monte-Carlo correctement construit, quelle que soit la technique
employée. C'est un \textbf{diagnostic}.

Une pente qui s'en écarterait signalerait un estimateur mal formé --- typiquement
une erreur type calculée sur un échantillon antithétique aplati, où les tirages
ne sont pas indépendants.

Le gain des techniques de réduction de variance ne se lit donc pas dans la
pente --- les quatre droites de la figure 10 sont parallèles --- mais dans leur
\emph{hauteur}. Aucune technique ne change le régime de convergence : elles
translatent la droite vers le bas.
\end{encadre}

\subsection{La pente n'est jamais imposée}

Une version antérieure du code portait le commentaire « la pente doit valoir
0.5 », formulation malheureuse : elle laisse croire que la valeur est
contrainte. Elle ne l'est pas. La pente est le résultat brut d'une régression,
elle n'est réinjectée dans aucun calcul, et le script affiche désormais
l'\textbf{écart à 0,50} pour que le lecteur juge lui-même :

\begin{center}
\small
\begin{tabular}{lrr}
\toprule
Méthode (asiatique) & pente mesurée & écart à 0,50 \\
\midrule
Monte-Carlo seul       & 0,501 & $+$0,001 \\
$+$ antithétique       & 0,505 & $+$0,005 \\
$+$ variable de contrôle & 0,508 & $+$0,008 \\
$+$ les deux           & 0,507 & $+$0,007 \\
\bottomrule
\end{tabular}
\end{center}

Les écarts sont de l'ordre du centième, ce qui est le bruit attendu d'une
régression sur quatre points bruités. Si l'un d'eux valait 0,3 ou 0,7, il
faudrait chercher l'erreur.

\begin{choix}
\textbf{Choix : titre de la figure 10.} L'ancien titre, « les quatre variantes
ont la même pente », annonçait la conclusion avant de la montrer. Le titre est
désormais simplement \textbf{« Erreur type selon $M$ »}, et les pentes mesurées
figurent dans la légende, avec trois décimales.
\end{choix}

\section{Le nombre de trajectoires nécessaire : mesuré, non extrapolé}

\ou{tableau 8, colonne « $M$ requis », et tableau 10.}

\subsection{Ce qui a changé, et pourquoi}

L'ancienne version extrapolait : elle mesurait l'erreur type sur un échantillon,
puis appliquait la loi en $1/\sqrt{M}$ pour calculer le $M$ qui \emph{aurait}
donné le centime. Le résultat était correct, mais indirect : on annonçait un
nombre de trajectoires qu'on n'avait jamais simulé, et le calcul demandait au
lecteur d'admettre une formule.

La nouvelle version \textbf{mesure}. C'est plus simple à expliquer et ça ne
repose sur rien.

\subsection{Comment la mesure procède}

\begin{enumerate}\setlength\itemsep{1pt}
\item On tire un paquet de trajectoires.
\item On met à jour trois compteurs : nombre de tirages, somme des gains, somme
      des carrés. Ces trois nombres suffisent à recalculer la moyenne et
      l'écart-type de \emph{tout} ce qui a été tiré depuis le début.
\item Dès que la demi-largeur de l'intervalle de confiance cumulé passe sous le
      centime, on s'arrête et on rapporte le nombre de trajectoires simulées.
\end{enumerate}

Les compteurs évitent de garder les trajectoires en mémoire : seuls trois
nombres passent d'un paquet au suivant. C'est ce qui permet d'atteindre neuf
millions de trajectoires sur 252 pas sans jamais dépasser quelques centaines de
mégaoctets.

\begin{choix}
\textbf{Choix : la taille des paquets.} 500 trajectoires au début, puis un
huitième de ce qui est déjà tiré, plafonné à 200\,000. Petits paquets au départ
pour bien localiser les réponses faibles (le contrôle géométrique s'arrête à
2\,000) ; gros paquets ensuite pour ne pas perdre de temps en appels quand la
réponse se compte en millions. Ce réglage ne change pas le résultat, seulement
la finesse avec laquelle on le localise : le $M$ rapporté est le premier
multiple de la taille de paquet qui atteint la cible.
\end{choix}

\begin{choix}
\textbf{Choix : la cible de $\pm 0{,}01$.} Un centime, c'est le pas de
cotation. Si l'incertitude vaut $\pm 0{,}09$, on ne peut pas annoncer
« 10,45 » : on ne sait pas si c'est 10,36 ou 10,55. Viser mieux que le centime
n'aurait pas de sens non plus, l'écart entre prix acheteur et prix vendeur sur
le marché étant bien plus large.
\end{choix}

\subsection{Résultats mesurés}

Toutes les valeurs ci-dessous sont le nombre de trajectoires \emph{effectivement
simulées} pour que l'intervalle de confiance passe sous $\pm 0{,}01$.

\begin{center}
\small
\begin{tabular}{lrrrrr}
\toprule
Méthode & pente & écart à 0,50 & $M$ mesuré & prix obtenu & IC95 \\
\midrule
\multicolumn{6}{l}{\textbf{Call européen} ($n = 1$) --- référence exacte 10,4506}\\
Monte-Carlo seul        & 0,497 & $-$0,003 & 8\,422\,106 & 10,4487 & 0,00994 \\
$+$ antithétique        & 0,502 & $+$0,002 & 4\,222\,106 & 10,4552 & 0,00992 \\
$+$ contrôle            & 0,499 & $-$0,001 & 1\,281\,664 & 10,4545 & 0,00971 \\
$+$ les deux            & 0,502 & $+$0,002 & 2\,422\,106 & 10,4533 & 0,00987 \\
\midrule
\multicolumn{6}{l}{\textbf{Asiatique arithmétique} ($n = 252$) --- référence 5,7822 $\pm$ 0,0004}\\
Monte-Carlo seul        & 0,501 & $+$0,001 & 2\,622\,106 & 5,7812 & 0,00967 \\
$+$ antithétique        & 0,505 & $+$0,005 & 1\,281\,664 & 5,7729 & 0,00961 \\
$+$ contrôle géométrique & 0,508 & $+$0,008 & \textbf{2\,000} & 5,7764 & 0,00964 \\
$+$ les deux            & 0,507 & $+$0,007 & 2\,000 & 5,7739 & 0,00944 \\
\midrule
\multicolumn{6}{l}{\textbf{Barrière up-and-out} $B = 120$ --- référence 1,3271 $\pm$ 0,0077}\\
Monte-Carlo seul        & 0,499 & $-$0,001 & 444\,024 & 1,3246 & 0,00999 \\
$+$ antithétique        & 0,498 & $-$0,002 & 394\,688 & 1,3270 & 0,00995 \\
$+$ contrôle vanille    & 0,499 & $-$0,001 & 444\,024 & 1,3247 & 0,00999 \\
\midrule
\multicolumn{6}{l}{\textbf{Barrière up-and-in} $B = 120$ --- référence 9,1235 $\pm$ 0,0077}\\
Monte-Carlo seul        & 0,500 & $-$0,000 & 8\,822\,106 & 9,1217 & 0,00999 \\
$+$ contrôle vanille    & 0,499 & $-$0,001 & 444\,024 & 9,1258 & 0,00999 \\
\bottomrule
\end{tabular}
\end{center}

Trois lectures.

\textbf{Le contrôle géométrique divise l'effort par 1311} sur l'asiatique :
2\,000 trajectoires au lieu de 2,6 millions. C'est le résultat le plus
spectaculaire du projet, et il vient entièrement d'une corrélation de 0,9996
entre les deux payoffs.

\textbf{Sur la up-and-out, le contrôle ne sert à rien} : 444\,024 trajectoires
avec ou sans lui, au chiffre près. La corrélation avec la vanille est nulle.

\textbf{Sur la up-and-in, le même contrôle divise l'effort par 20} : 444\,024
au lieu de 8,8 millions. Et l'on retrouve exactement le nombre de la up-and-out,
ce qui n'est pas un hasard : contrôler la « in » par la vanille revient
algébriquement à estimer « vanille moins out », donc à hériter de la variance
de la out.

\begin{encadre}
\textbf{Cohérence avec l'ancienne extrapolation.} Sur le call européen sans
réduction de variance, l'extrapolation annonçait 8\,331\,880 trajectoires ; la
mesure directe en trouve 8\,422\,106, soit 1 \% d'écart. La loi en $1/\sqrt{M}$
était donc bien vérifiée --- mais il n'y a plus besoin de l'invoquer.
\end{encadre}

\section{Les variables de contrôle : pourquoi ces choix}

\ou{section 7.5.2 et tableau 10.}

\subsection{Le principe, sans formule}

On cherche à valoriser un produit A dont on ne connaît pas le prix. On repère
un produit B qui \emph{ressemble} à A et dont on connaît, lui, le prix exact.
On simule les deux avec les mêmes tirages. On regarde de combien la simulation
se trompe sur B --- puisqu'on connaît la vérité pour B --- et on suppose qu'elle
se trompe de façon comparable sur A. On retranche cette erreur.

\begin{encadre}
\textbf{Ce qui décide de tout : la corrélation.} Si la corrélation entre les
deux payoffs vaut $\rho$, la variance résiduelle est réduite d'un facteur
$1-\rho^2$. Une corrélation de 0,99 divise la variance par 50 ; une corrélation
nulle n'apporte rien. \textbf{Choisir un contrôle, c'est choisir une
corrélation.}
\end{encadre}

\subsection{Européenne : le sous-jacent lui-même}

Contrôle : le prix terminal $S_T$ actualisé. Son espérance est connue sans le
moindre calcul : sous la probabilité risque-neutre, le prix actualisé est une
martingale, donc son espérance vaut le spot d'aujourd'hui corrigé du dividende.
Corrélation mesurée avec un call à la monnaie : \textbf{environ 0,90}. Gain
modeste mais gratuit.

\subsection{Asiatique : la version géométrique}

Contrôle : l'asiatique géométrique, qui possède une formule fermée. C'est le
choix classique du domaine, introduit par Kemna et Vorst en 1990.

Intuition : les deux options portent sur \emph{la même trajectoire} et
regardent \emph{les mêmes 252 prix}. Seule la façon de les moyenner diffère.
Elles ne peuvent quasiment pas se désolidariser.

\begin{center}
\small
\begin{tabular}{lr}
\toprule
Corrélation mesurée (arithmétique, géométrique) & \textbf{0,99962} \\
Réduction de variance $1-\rho^2$ & 0,00075 \\
Gain sur la variance & facteur \textbf{1330} \\
Gain sur l'erreur type & facteur 36,5 \\
$\beta$ optimal estimé & 1,0350 \\
\bottomrule
\end{tabular}
\end{center}

Le facteur 1330 théorique et le facteur 1311 mesuré au tableau précédent sont
le même nombre, à la résolution des paquets près.

\subsection{Barrière : la vanille de mêmes caractéristiques}

Contrôle : la vanille de même strike et même maturité, dont Black--Scholes donne
le prix exact. L'intuition est que la barrière \emph{est} la vanille, amputée
des trajectoires qui franchissent le seuil.

\begin{center}
\small
\begin{tabular}{lcc}
\toprule
& Corrélation avec la vanille & Gain mesuré \\
\midrule
up-and-\textbf{in}  & $+0{,}974$ & $\times 20$ \\
up-and-\textbf{out} & $-0{,}010$ & \textbf{aucun} \\
\bottomrule
\end{tabular}
\end{center}

\textbf{Pourquoi.} La up-and-in ne paie que si la barrière est touchée, donc
dans les scénarios franchement haussiers --- précisément ceux où la vanille paie
le plus. Les deux montent ensemble. La up-and-out, au contraire, ne paie que
dans une bande étroite : finir au-dessus du strike \emph{sans jamais} dépasser
la barrière. Les trajectoires les plus rentables pour la vanille sont
exactement celles qui la tuent. La corrélation s'annule mécaniquement.

\textbf{Et il n'existe aucun détour.} Contrôler la « in » par la vanille redonne
« vanille moins out » : la variance résiduelle reste celle de la out. Ce produit
est la partie irréductible du problème. La littérature apporte une réponse,
mais d'une autre nature : un conditionnement sur la survie, et non une variable
de contrôle (Glasserman et Staum, 2001).

\subsection{Pourquoi les antithétiques échouent aussi sur la barrière}

Les variables antithétiques supposent un payoff monotone en l'aléa. Une
barrière ne l'est pas : un choc très positif désactive l'option (gain nul), un
choc très négatif la laisse hors de la monnaie (gain nul également). Les deux
extrêmes donnent zéro, la corrélation négative recherchée disparaît. Le tableau
mesuré le confirme : gain de 1,1 seulement.

\begin{choix}
\textbf{Choix : $\beta$ estimé sur le même échantillon que le prix.} La théorie
recommande un échantillon pilote séparé, car estimer $\beta$ sur les données
qu'on corrige introduit un biais d'ordre $1/M$ (Glasserman, \S 4.1.3). À
$M \geq 10^5$ ce biais est négligeable devant l'erreur type, et la simplicité a
été préférée. C'est un choix, pas une omission.
\end{choix}

\section{Récapitulatif des choix assumés}

\begin{center}
\small
\renewcommand{\arraystretch}{1.15}
\begin{tabular}{p{4.1cm}p{3.0cm}p{7.6cm}}
\toprule
Décision & Valeur retenue & Pourquoi \\
\midrule
Fixings des exotiques & $n = 252$ & jours de bourse ; doit être cité avec chaque prix \\
Cible de précision & $\pm$ 0,01 & le pas de cotation ; en dessous, calcul gaspillé \\
Niveau de confiance & 95 \% ($\times 1{,}96$) & convention usuelle \\
Lots pour la référence & 10 $\times$ 100\,000 & mémoire ; 5 lots donnent une incertitude trop bruitée \\
Graine des références & 999 & différente de 12345 pour éviter le recouvrement d'échantillons \\
Graine des mesures & 12345 & fixée pour la reproductibilité \\
Taille des paquets & 500, puis $M/8$, max 200\,000 & compromis résolution / vitesse \\
$\beta$ du contrôle & estimé sur l'échantillon & biais en $1/M$ négligeable ici \\
Antithétiques + contrôle & combinés & se composent, une fois $\beta$ estimé sur les moyennes de paires (voir ci-dessous) \\
\bottomrule
\end{tabular}
\end{center}

\paragraph{Rectificatif : antithétiques et variable de contrôle se composent.}
Une version antérieure de cette note indiquait que combiner les deux techniques
faisait moins bien que le contrôle seul. C'était la conséquence d'un défaut
d'implémentation, depuis corrigé : le coefficient $\beta$ était estimé sur
l'échantillon \emph{aplati} au lieu des moyennes de paires. Or l'antithétique
annule la composante impaire de l'aléa ; le contrôle aplati visait précisément
cette composante, donc ne trouvait plus rien à corriger. Régressé sur les
moyennes de paires, le même contrôle devient une fonction paire de l'aléa et
attaque le complément. Sur le call européen à la monnaie, $M=400\,000$ :
l'erreur type passe de $0{,}01237$ à $0{,}00433$, la corrélation de $0{,}9245$ à
$0{,}9644$, et $\beta$ de $0{,}674$ à $2{,}508$. Le nombre de trajectoires mesuré
pour coter au centime passe de $2\,422\,106$ à $311\,854$. Le prix, lui, n'était
pas affecté : l'estimateur à variable de contrôle est sans biais quelle que soit
la valeur de $\beta$. Voir le rapport, section 7.5, et
\texttt{CORRECTIONS.md} §9.

\section{Sources}

\newcommand{\refitem}[3]{\textbf{#1}\\ #2 \\ {\footnotesize\url{#3}}\\[4pt]}

\refitem{Kemna, A.G.Z. et Vorst, A.C.F. (1990)}{%
« A pricing method for options based on average asset values », \emph{Journal
of Banking \& Finance}, 14(1), 113--129. --- \textsl{L'article fondateur du
contrôle géométrique pour l'asiatique arithmétique.}}{%
https://doi.org/10.1016/0378-4266(90)90039-5}

\refitem{Boyle, P. (1977)}{%
« Options: A Monte Carlo Approach », \emph{Journal of Financial Economics},
4(3), 323--338. --- \textsl{Introduction du Monte-Carlo et des variables de
contrôle en valorisation d'options.}}{%
https://doi.org/10.1016/0304-405X(77)90005-8}

\refitem{Boyle, P., Broadie, M. et Glasserman, P. (1997)}{%
« Monte Carlo methods for security pricing », \emph{Journal of Economic
Dynamics and Control}, 21(8--9), 1267--1321. --- \textsl{Synthèse de référence :
coefficient optimal, gain en $1-\rho^2$, efficacité intégrant le temps de
calcul. Version libre d'accès :}}{%
https://www.cs.hunter.cuny.edu/~felisav/StochasticProcesses/Course_Material_files/BBG_jedc96.pdf}

\refitem{Glasserman, P. (2004)}{%
\emph{Monte Carlo Methods in Financial Engineering}, Springer, chapitre 4.
--- \textsl{Variables de contrôle (\S 4.1), antithétiques (\S 4.2), biais
d'estimation de $\beta$ (\S 4.1.3).}}{%
https://doi.org/10.1007/978-0-387-21617-1}

\refitem{Glasserman, P. et Staum, J. (2001)}{%
« Conditioning on One-Step Survival for Barrier Option Simulations »,
\emph{Operations Research}, 49(6), 923--937. --- \textsl{La technique qui
fonctionne sur la barrière knock-out, là où la variable de contrôle
échoue.}}{%
https://doi.org/10.1287/opre.49.6.923.10018}

\refitem{Turnbull, S. et Wakeman, L. (1991)}{%
« A Quick Algorithm for Pricing European Average Options », \emph{Journal of
Financial and Quantitative Analysis}, 26(3), 377--389. --- \textsl{Approximation
analytique de l'asiatique arithmétique : alternative au contrôle, utilisable
comme seconde référence.}}{%
https://doi.org/10.2307/2331214}

\end{document}
